
% =========================================================
% II. RELATED WORK
% =========================================================
\section{Related Work}
\label{sec:related_work}

\subsection{Retrieval-Augmented Generation}

Retrieval-Augmented Generation (RAG) combines a language model with an external
retrieval component to provide relevant evidence prior to answer generation
\cite{lewis2020retrieval,gao2024ragsurvey}. Dense retrieval methods such as DPR
\cite{karpukhin2020dense} improve semantic document retrieval by representing
queries and passages in a shared embedding space. Subsequent work has explored
the aggregation of multiple retrieved passages, adaptive retrieval and
self-reflection, and hierarchical retrieval structures
\cite{izacard2021fid,asai2024selfrag,sarthi2024raptor}. However, most
conventional RAG systems still provide retrieved evidence primarily as textual
context rather than as an explicit reasoning structure.

This paradigm is effective for questions that can be answered using one or a
small number of directly relevant passages. Its effectiveness is more limited,
however, for tasks that require information to be connected across multiple
reasoning steps.

\subsection{Multi-hop Question Answering}

Multi-hop Question Answering requires a system to integrate information from
multiple pieces of evidence in order to derive a final answer. Benchmarks such
as HotpotQA \cite{yang2018hotpotqa} and MuSiQue \cite{trivedi2022musique}
demonstrate that retrieving individually relevant documents is not sufficient;
a system must also identify the reasoning chain that connects the supporting
evidence.

This requirement is particularly important in the legal domain, where a legal
conclusion may depend on multiple conditions, exceptions, and intermediate
consequences. Explicitly representing such reasoning structures may therefore
help reduce unsupported reasoning transitions and improve the traceability of
generated conclusions.

\subsection{Graph-based and Causal RAG}

Recent research has explored graph-based representations as a means of
organizing knowledge and supporting multi-hop retrieval. GraphRAG leverages
graph structures to expand retrieval and aggregate information across related
entities and passages \cite{edge2024graphrag}. Think-on-Graph further
demonstrates that explicit traversal over a knowledge graph can support
multi-step reasoning and provide traceable reasoning paths
\cite{sun2024thinkongraph}.

More closely related to our work, CausalRAG incorporates causal relations into
retrieval and reasoning rather than relying solely on semantic similarity
\cite{wang2025causalrag}. CC-RAG similarly constructs causal chains to support
structured multi-hop inference over domain-specific corpora
\cite{parekh2025ccrag}. This line of work suggests that explicit relations
among events can provide useful structural signals for multi-hop reasoning.

In the legal setting considered in this study, however, these relations are
interpreted as directed condition--consequence relations derived from statutory
text. They should therefore be understood as a structured representation of
legal dependencies rather than as causal effects estimated from observational
data.

\subsection{Counterfactual Reasoning and Legal QA}

Counterfactual reasoning provides a mechanism for examining whether a
conclusion would remain valid if an intermediate condition or event were
modified or removed. In multi-hop question answering, prior work on
counterfactual questions has shown that a model may arrive at the correct
answer while relying on an inappropriate, incomplete, or disconnected
reasoning chain \cite{guo2023counterfactual}.

In the legal domain, benchmarks such as LexGLUE, CaseHOLD, and LegalBench, as
well as domain-specific models such as LEGAL-BERT, have highlighted both the
distinctive characteristics of legal language and the continuing difficulty of
structured legal reasoning
\cite{chalkidis2022lexglue,zheng2021casehold,guha2023legalbench,chalkidis2020legalbert}.
These limitations motivate approaches that combine evidence retrieval with more
explicit representations of legal relations and reasoning structure.

\subsection{Three-Valued and Linear-Algebraic Logic Programming}

Three-valued logic distinguishes true, false, and an undefined or unknown state
\cite{kleene1952metamathematics}. Fitting's Kripke--Kleene semantics provides a
fixed-point account for logic programs under partial information
\cite{fitting1985kripke}. More recent work encodes two- and three-valued logic
program semantics in vector spaces \cite{sato2020threevalued}, while matrix and
third-order tensor representations support definite, disjunctive, and normal
logic programs \cite{sakama2021tensor}. These studies motivate our algebraic
formulation. Our runtime remains a symbolic sparse-rule engine: it uses AND
within one rule and OR across independently activated rules with the same
consequent rather than executing tensor multiplication.

\subsection{Research Gap}

Conventional RAG primarily focuses on retrieving semantically relevant
evidence, whereas graph-based approaches emphasize structural connections
among entities, events, or passages. Nevertheless, relatively limited research
has investigated legal QA frameworks that jointly incorporate the following
capabilities:

\begin{itemize}
    \item representing legal rules as directed condition--consequence relations;
    \item retrieving and ranking multi-hop reasoning paths;
    \item assessing the necessity of intermediate events through counterfactual
    analysis; and
    \item using verification outcomes to refine the evidence provided to the
    answer-generation model.
\end{itemize}

This work addresses this gap by integrating causal graph retrieval, multi-hop
path reasoning, and counterfactual verification within a unified RAG framework
for Vietnamese multi-hop legal question answering.
